%%%%%%%%%%%%%%%%%%%%%%%%%%%%%%%%%%%%%%%%%%%%%%%%%%%%%%%%%%%%%%%%%%%%%%%%%%%%%%%%
% DT4TM -- Endabgabe (Code-Dokumentation + Methodik FEM -> MOR -> Zwilling)
% Stand: 26.09.2026 (Version 1). Lesefassung: project_course_final.md
% Gleiche Vorlage wie project_course_article.tex (ieeeconf, 10pt, 2-spaltig).
% Figures/rom_vs_fem.pdf wird erzeugt aus: FEM-Transiente vs. ROM (siehe Abschn. VI).
%%%%%%%%%%%%%%%%%%%%%%%%%%%%%%%%%%%%%%%%%%%%%%%%%%%%%%%%%%%%%%%%%%%%%%%%%%%%%%%%

\documentclass[a4paper, 10pt, conference]{ieeeconf}
\IEEEoverridecommandlockouts
\overrideIEEEmargins

\usepackage[utf8]{inputenc}
\usepackage[T1]{fontenc}
\usepackage[ngerman]{babel}
\usepackage{graphicx}
\usepackage{amsmath}
\usepackage[left=1.57cm,right=1.57cm,top=2.54cm,bottom=2.54cm]{geometry}
\usepackage{amssymb}
\usepackage{enumitem}
\usepackage{amsthm}
\usepackage{booktabs}
\usepackage{tikz}
\usetikzlibrary{arrows.meta, positioning, fit, backgrounds}
\usepackage{url}
\usepackage{tabularx}
\usepackage[hidelinks]{hyperref}

\theoremstyle{remark}
\newtheorem{remark}{Anmerkung}

\makeatletter
\def\@IEEEprocessthesectionargument#1{%
\@ifmtarg{#1}{%
\@IEEEappendixsavesection*{Anhang \thesectiondis}%
\addcontentsline{toc}{section}{Anhang \thesection}}{%
\@IEEEappendixsavesection*{Anhang \thesectiondis \\* #1}%
\addcontentsline{toc}{section}{Anhang \thesection: #1}}}
\makeatother

% Kasten "Kernidee": fasst jeden Methodenschritt in 2-3 Sätzen für Leser ohne
% Simulationshintergrund zusammen. Nur Bordmittel (\fbox), kein Zusatzpaket.
\newcommand{\kernidee}[1]{\par\smallskip\noindent
  \fbox{\parbox{\dimexpr\columnwidth-2\fboxsep-2\fboxrule\relax}{%
  \small\textbf{Kernidee.}~#1}}\par\smallskip}

\title{\LARGE \bf
Vom Feldmodell zum Echtzeit-Zwilling:\\
Modellordnungsreduktion für den thermischen Digital Twin\\
eines elektrodynamischen Levitators
}
\author{Dang Minh Hoang, Marina Borchert, Mohamed Aziz El Majid
\thanks{Project Course Digital Twin, Summer Term 2026. Supervisors: Hr. Prof. Dr. rer. nat. Dirk Hartmann, Fr. Dr.-Ing. Melina Merkel. Stand: 26.09.2026}%
}

\begin{document}
\maketitle
\thispagestyle{plain}
\pagestyle{plain}

%%%%%%%%%%%%%%%%%%%%%%%%%%%%%%%%%%%%%%%%%%%%%%%%%%%%%%%%%%%%%%%%%%%%%%%%%%%%%%%%
\begin{abstract}
Ein digitaler Zwilling soll die Temperatur von Scheibe und Spulen eines
TEAM-28-ähnlichen Levitators vorhersagen, und zwar laufend, aus dem
gemessenen Spulenstrom, auf einem Smartphone. Dieser Bericht erklärt den
Weg dorthin als Kette aus vier Schritten: (1) Physik als partielle
Differentialgleichungen, (2) deren Lösung mit der Finite-Elemente-Methode
(FEM) als großes lineares Gleichungssystem, (3) die Reduktion dieses
Systems auf wenige Zustandsgrößen (\emph{Model Order Reduction}, MOR) und
(4) die Kalibrierung der reduzierten Parameter an Messdaten. Kern des
Berichts ist die Begründung von Schritt~3: Aus $8700$ komplexen
EM-Unbekannten und $697$ thermischen Unbekannten wird ein Modell mit acht
thermischen Zuständen, dessen Zeitschritt $\approx3\,\mu\mathrm s$ dauert.
Eine Eigenwertanalyse der Wärme-FEM zeigt, warum das zulässig ist: Die
langsamste Zeitkonstante ($204\,\mathrm s$) ist $47$-mal größer als die
nächste, eine einzige Mode gibt die volle FEM-Transiente daher mit
$\le0{,}36\,\mathrm K$ Fehler wieder. Zugleich dokumentiert der Bericht den
Code so, dass seine Funktion ohne Lektüre des Quelltexts nachvollziehbar
ist, und fasst Endergebnisse, offene Punkte und Lessons Learned zusammen.
\end{abstract}

%%%%%%%%%%%%%%%%%%%%%%%%%%%%%%%%%%%%%%%%%%%%%%%%%%%%%%%%%%%%%%%%%%%%%%%%%%%%%%%%
\section{Einleitung: Was soll der Zwilling leisten?}
\label{sec:intro}

Der Versuchsstand ist ein elektrodynamischer Levitator nach dem
TEAM-Benchmark~28~\cite{team28}: Zwei Wechselstromspulen ($50\,\mathrm{Hz}$)
induzieren in einer Aluminiumscheibe Wirbelströme, die Scheibe schwebt.
Dieselben Ströme erzeugen Wärme in Scheibe, Spulen und Eisen.
Abbildung~\ref{fig:geometry} zeigt den Aufbau, Tabelle~\ref{tab:geometry}
die vermessenen Maße.

Die Temperatur in der schwebenden Scheibe oder im Inneren der Wicklung
lässt sich im Betrieb kaum messen: Ein Thermoelement an der Scheibe stört
die Levitation, und blankes Aluminium ist für die IR-Kamera wegen seiner
unbekannten Emissivität unzuverlässig. Genau das ist der Anwendungsfall
\emph{virtueller Sensor}, den Hartmann und Van der Auweraer für den
Anlauf großer Elektromotoren beschreiben~\cite{hartmann2025golden}: Eine
nicht messbare Temperatur wird aus einer messbaren Größe (dort
Statorsensoren, hier der Spulenstrom) über ein physikalisches Modell
berechnet. Daraus folgen drei Anforderungen:
\begin{enumerate}[label=(\roman*), leftmargin=*, nosep]
  \item \textbf{Echtzeit:} Das Modell muss mit dem Messtakt ($1\,\mathrm{Hz}$)
        mitlaufen, auch auf einem Smartphone im Browser.
  \item \textbf{Physikalische Treue:} Es muss aus Maxwell- und
        Wärmeleitungsgleichung abgeleitet sein, nicht bloß an Daten
        angepasst.
  \item \textbf{Kalibrierbarkeit:} Die wenigen unsicheren Parameter
        (Wärmeübergang) müssen sich an Messungen nachstellen lassen.
\end{enumerate}
(i) und (ii) widersprechen sich auf den ersten Blick: Genaue Modelle sind
groß und langsam, schnelle Modelle ungenau. Die Methode, die beide
Anforderungen verbindet, ist die Modellordnungsreduktion, die Hartmann et
al.\ als Schlüsseltechnologie für digitale Zwillinge
bezeichnen~\cite{hartmann2018mor}. Dieser Bericht erklärt sie am
konkreten Beispiel. Abbildung~\ref{fig:chain} dient dabei als Leitfaden:
Die Abschnitte~II--VI behandeln je einen Kasten.

\begin{figure}[t]
\centering
\begin{tikzpicture}[x=0.065cm, y=0.065cm,
    ironfill/.style={fill=gray!45},
    cufill/.style={fill=orange!35},
    alfill/.style={fill=cyan!25},
    lbl/.style={font=\scriptsize, align=center}]
  \draw[dash dot] (0,15) -- (0,-58);
  \node[font=\tiny, right] at (2,16) {$r=0$};
  \draw[-{Stealth[length=4pt]}] (0,-58) -- (109,-58) node[below, font=\tiny] {$r$ [mm]};
  \draw (25.9,-58) -- (25.9,-56.5); \node[font=\tiny] at (25.9,-63) {25{,}9};
  \draw (61.9,-58) -- (61.9,-56.5); \node[font=\tiny] at (61.9,-63) {61{,}9};
  \draw (79.9,-58) -- (79.9,-56.5); \node[font=\tiny] at (79.9,-63) {79{,}9};
  \draw (102.9,-58) -- (102.9,-56.5); \node[font=\tiny] at (102.9,-63) {102{,}9};
  \draw[ironfill] (0,0) rectangle (25.9,-53);
  \node[lbl] at (13,-26.5) {Eisenkern\\$\mu_r{=}1000$};
  \draw[cufill] (27.9,0) rectangle (61.9,-53);
  \node[lbl] at (44.9,-26.5) {Innenspule\\$N{=}1000$};
  \draw[ironfill] (64.9,0) rectangle (79.9,-53);
  \node[lbl, rotate=90] at (72.4,-26.5) {Eisenring};
  \draw[cufill] (82.9,0) rectangle (102.9,-53);
  \node[lbl] at (92.9,-26.5) {Außen-\\spule\\$N{=}500$};
  \draw[alfill] (0,6) rectangle (80,9);
  \node[lbl] at (40,7.5) {Al-Scheibe, $R{=}80$, $d{=}3$\,[mm]};
  \draw[<->, dashed] (85,0) -- (85,6);
  \node[font=\tiny] at (91,3) {$z_{\mathrm{gap}}$};
\end{tikzpicture}
\caption{Halber Querschnitt des Geräts (rotationssymmetrisch um $r=0$),
identisch zur Eingabe in \texttt{params.yaml}. Weil sich um die Achse
nichts ändert, genügt es, diese 2D-Ebene $(r,z)$ zu rechnen.}
\label{fig:geometry}
\end{figure}

\begin{table}[t]
\centering\footnotesize
\caption{Vermessene Geometrie und Materialwerte (\texttt{params.yaml}).}
\label{tab:geometry}
\begin{tabular}{@{}lrrrr@{}}
\toprule
Bauteil & $r$\,[mm] & $\mu_r$ & $\sigma$\,[MS/m] & $N$ \\
\midrule
Eisenkern & 0{,}0--25{,}9 & 1000$^\ast$ & 1{,}0$^\ast$ & -- \\
Innenspule (Cu) & 27{,}9--61{,}9 & 1 & 59{,}6 & 1000 \\
Eisenring & 64{,}9--79{,}9 & 1000$^\ast$ & 1{,}0$^\ast$ & -- \\
Außenspule (Cu) & 82{,}9--102{,}9 & 1 & 59{,}6 & 500 \\
Al-Scheibe ($d=3\,$mm) & 0{,}0--80{,}0 & 1 & 34{,}0 & -- \\
\bottomrule
\end{tabular}

\vspace{2pt}
{\footnotesize Höhe von Kern, Spulen und Ring: $53\,\mathrm{mm}$.
Betriebspunkt $I_{\mathrm{rms}}=5\,\mathrm A$, $f=50\,\mathrm{Hz}$;
Kalibrierpunkt $7{,}8\,\mathrm A$. $^\ast$Literaturwerte, nicht gemessen
(Abschnitt~\ref{sec:results}).}
\end{table}

\begin{figure*}[t]
\centering
\resizebox{\textwidth}{!}{%
\begin{tikzpicture}[
  node distance=4mm and 7mm,
  step/.style={draw, rounded corners, align=center, text width=3.55cm,
               minimum height=2.1cm, font=\footnotesize, fill=white},
  arr/.style={-Stealth, thick},
  lab/.style={font=\scriptsize\itshape, align=center}
]
\node[step] (s1) {\textbf{1. Physik}\\[2pt]
  Maxwell (Phasor) + Fourier\\ PDE in $(r,z)$\\[2pt]
  \scriptsize unendlich viele Unbekannte\\ \texttt{params.yaml}, \texttt{config.py}};
\node[step, right=of s1] (s2) {\textbf{2. FEM}\\[2pt]
  $\mathbf K\,\mathbf x=\mathbf f$\\
  EM: $8700$ (komplex), Wärme: $697$\\[2pt]
  \scriptsize $0{,}44\,$s pro Betriebspunkt\\ \texttt{em\_solver.py}, \texttt{thermal\_solver.py}};
\node[step, right=of s2] (s3) {\textbf{3. MOR}\\[2pt]
  $\mathbf x\approx \mathbf V\,\boldsymbol\beta(t)$\\
  $8$ thermische Zustände\\[2pt]
  \scriptsize $\approx3\,\mu$s pro Zeitschritt\\ \texttt{rom.py}, \texttt{twin\_core.py}};
\node[step, right=of s3] (s4) {\textbf{4. Zwilling}\\[2pt]
  Messstrom $I(t)$ rein,\\ $T(r,z,t)$ + 3D/AR raus;\\ Parameter an Messung fitten\\[2pt]
  \scriptsize \texttt{data\_io.py}, \texttt{refit\_hA.py},\\ \texttt{build\_twin\_html\_fem.py}};
\draw[arr] (s1) -- (s2);
\draw[arr] (s2) -- node[above, lab]{einmal} (s3);
\draw[arr] (s3) -- node[above, lab]{laufend} (s4);
\begin{scope}[on background layer]
\node[fill=gray!12, rounded corners, fit=(s1)(s2), inner sep=6pt,
      label={[font=\footnotesize\bfseries]above:OFFLINE (Laptop, einmal pro Geometrie)}] {};
\node[fill=cyan!10, rounded corners, fit=(s3)(s4), inner sep=6pt,
      label={[font=\footnotesize\bfseries]above:ONLINE (Browser/Smartphone, jeder Zeitschritt)}] {};
\end{scope}
\draw[arr, dashed] (s4.south) -- ++(0,-0.55) -| node[pos=0.25, below, lab]{Kalibrierung: $hA$, $\tau$ aus Messdaten} (s3.south);
\end{tikzpicture}%
}
\caption{Leitfaden des Berichts und zugleich Architektur des Codes. Die
teure Rechnung (grau) läuft einmal; im Betrieb (blau) läuft nur das
reduzierte Modell. Unter jedem Schritt stehen die Dateien, die ihn
umsetzen (Details in Tabelle~\ref{tab:modules}).}
\label{fig:chain}
\end{figure*}

%%%%%%%%%%%%%%%%%%%%%%%%%%%%%%%%%%%%%%%%%%%%%%%%%%%%%%%%%%%%%%%%%%%%%%%%%%%%%%%%
\section{Schritt 1: Vom Phänomen zur Gleichung}
\label{sec:physics}

Die Physik ist eine Kette aus drei Gliedern: \emph{Strom $\to$
Magnetfeld $\to$ Wirbelstrom $\to$ Wärme $\to$ Temperatur}. Jedes Glied
ist eine Gleichung.

\textbf{Magnetfeld.} Wegen der Rotationssymmetrie hat das magnetische
Vektorpotential nur eine Komponente $A_\varphi(r,z)$. Da der Strom
sinusförmig ist, wird nicht der Zeitverlauf gerechnet, sondern Amplitude
und Phase als komplexe Zahl, der \emph{Phasor}~\cite{biro1989vectorpotential}:
\begin{equation}
  -\nabla\!\cdot(\nu\nabla A_\varphi) + \frac{\nu}{r^2}A_\varphi
  + j\omega\sigma A_\varphi = J_s,\qquad \nu=\frac{1}{\mu_r\mu_0}.
  \label{eq:em}
\end{equation}
\emph{Lesehilfe:} Der erste Term beschreibt, wie sich das Feld im Raum
ausbreitet (Eisen mit großem $\mu_r$ bündelt es), der zweite ist der
geometrische Preis der Zylinderkoordinaten, der dritte ist die Gegenwehr
jedes Leiters: Das Wechselfeld induziert dort Wirbelströme, die dem Feld
entgegenwirken. $J_s$ ist die Stromdichte in den Spulen, die Quelle.

\textbf{Wärmequelle.} Aus dem Feld folgt die über eine Periode gemittelte
Verlustdichte im Leiter:
\begin{equation}
  q(r,z) = \tfrac12\,\sigma\,\omega^2\,|A_\varphi|^2 \quad[\mathrm{W/m^3}].
  \label{eq:q}
\end{equation}
Die Spulen selbst erwärmen sich ohmsch mit $P_{\mathrm{Cu}}=\tfrac12 I^2R$.

\textbf{Temperatur.} Die Wärme verteilt sich nach der Fourier-Gleichung
und wird an der Oberfläche durch Konvektion abgegeben:
\begin{equation}
  \rho c_p\frac{\partial T}{\partial t}=\nabla\!\cdot(k\nabla T)+q,
  \qquad -k\frac{\partial T}{\partial n}=h\,(T-T_\infty).
  \label{eq:heat}
\end{equation}

\kernidee{Zwei Zeitskalen: Das Feld schwingt mit $50\,\mathrm{Hz}$ ($20\,$ms),
die Temperatur ändert sich über Minuten. Das sind vier bis fünf
Größenordnungen. Die Thermik spürt nur den Mittelwert der Verluste,
deshalb wird das Feld einmal als Phasor gelöst und nie im
$50\,$Hz-Takt zeitlich integriert. Das ist die erste, noch ganz
physikalische Vereinfachung.}

%%%%%%%%%%%%%%%%%%%%%%%%%%%%%%%%%%%%%%%%%%%%%%%%%%%%%%%%%%%%%%%%%%%%%%%%%%%%%%%%
\section{Schritt 2: Wie die FEM die Gleichung löst}
\label{sec:fem}

Gl.~\eqref{eq:em} und \eqref{eq:heat} haben für diese Geometrie keine
geschlossene Lösung. Die FEM ersetzt die gesuchte Funktion durch endlich
viele Zahlen in drei Schritten:

\textbf{(a) Zerlegen.} Die Ebene $(r,z)$ wird in Dreiecke zerlegt
(Abb.~\ref{fig:fem}). Unbekannt ist nur noch der Wert in jedem
Eckpunkt (Knoten); dazwischen wird linear interpoliert (P1-Elemente). Jeder
Knoten $i$ besitzt eine „Zeltfunktion“ $v_i$, die an ihm $1$ und an allen
anderen Knoten $0$ ist.

\textbf{(b) Bilanz je Knoten (schwache Form).} Statt die Gleichung in
jedem Punkt zu erfüllen, verlangt man, dass sie \emph{gewichtet mit jeder
Zeltfunktion} im Mittel erfüllt ist. Für die Wärmeleitung heißt das
anschaulich: Um jeden Knoten herum muss die hineingeleitete, erzeugte und
abgegebene Wärme bilanziert sein. Mathematisch (stationär, mit dem
Volumenfaktor $2\pi r$ der Rotation):
\begin{equation}
  \int_\Omega k\nabla T\!\cdot\!\nabla v_i\,2\pi r\,dA
  +\oint_\Gamma hTv_i\,2\pi r\,ds
  =\int_\Omega q\,v_i\,2\pi r\,dA+\oint_\Gamma hT_\infty v_i\,2\pi r\,ds .
  \label{eq:weak}
\end{equation}

\textbf{(c) Gleichungssystem.} Setzt man $T=\sum_j T_jv_j$ ein, wird
aus \eqref{eq:weak} für alle Knoten zusammen
\begin{equation}
  \mathbf K\,\mathbf T=\mathbf f,
  \qquad\text{instationär:}\quad
  \mathbf M\,\dot{\mathbf T}+\mathbf K\,\mathbf T=\mathbf f .
  \label{eq:linsys}
\end{equation}
$\mathbf K$ (Leitung + Konvektion) und $\mathbf M$ (Wärmekapazität) sind
dünn besetzt, weil jeder Knoten nur mit seinen Nachbarn koppelt. Für das
EM-Problem entsteht dasselbe mit komplexem $\mathbf K$ (wegen
$j\omega\sigma$). Beide Systeme werden mit dem direkten Löser
\texttt{scipy.sparse.linalg.spsolve}~\cite{virtanen2020scipy} gelöst.

\begin{figure}[t]
\centering
\begin{tikzpicture}[scale=0.95]
  % kleines Dreiecksnetz
  \foreach \i in {0,...,4}{\foreach \j in {0,1,2}{
    \coordinate (p\i\j) at (\i*0.9,\j*0.8);}}
  \foreach \i [evaluate=\i as \ii using int(\i+1)] in {0,...,3}{
    \foreach \j [evaluate=\j as \jj using int(\j+1)] in {0,1}{
      \draw[gray] (p\i\j) -- (p\ii\j) -- (p\ii\jj) -- (p\i\jj) -- cycle;
      \draw[gray] (p\i\j) -- (p\ii\jj);}}
  % Zeltfunktion (Träger) um Knoten p21
  \fill[orange!35, opacity=0.8] (p11) -- (p10) -- (p20) -- (p31) -- (p32) -- (p22) -- cycle;
  \foreach \i in {0,...,4}{\foreach \j in {0,1,2}{\fill (p\i\j) circle (1.1pt);}}
  \fill[red] (p21) circle (2pt);
  \node[font=\scriptsize, above right] at (p21) {$i$};
  \draw[-Stealth] (-0.3,-0.35) -- (4.0,-0.35) node[right, font=\scriptsize] {$r$};
  \draw[-Stealth] (-0.3,-0.35) -- (-0.3,1.9) node[above, font=\scriptsize] {$z$};
  % 1D-Schnitt der Zeltfunktion
  \begin{scope}[xshift=4.6cm, yshift=0.2cm]
    \draw[-Stealth] (0,0) -- (2.6,0) node[right, font=\scriptsize] {$r$};
    \draw[thick, orange!80!black] (0,0) -- (0.4,0) -- (1.2,1.1) -- (2.0,0) -- (2.4,0);
    \node[font=\scriptsize] at (1.2,1.35) {$v_i$};
    \node[font=\scriptsize, below] at (1.2,0) {$r_i$};
    \draw[dashed] (1.2,0) -- (1.2,1.1);
  \end{scope}
\end{tikzpicture}
\caption{Idee der FEM: Die Temperatur ist nur in den Knoten unbekannt.
Die Zeltfunktion $v_i$ (orange: ihr Träger, rechts: 1D-Schnitt) koppelt
Knoten $i$ nur mit seinen direkten Nachbarn, daher ist $\mathbf K$ fast
überall null.}
\label{fig:fem}
\end{figure}

\textbf{Warum 2D genügt und das Netz grob sein darf.} Die Rotationssymmetrie
reduziert 3D auf 2D. Die Skintiefe in Aluminium bei $50\,\mathrm{Hz}$,
$\delta=\sqrt{2/(\omega\mu\sigma)}\approx12\,\mathrm{mm}$, ist viermal
größer als die Scheibendicke, der Wirbelstrom ist also über die Dicke fast
konstant; ein feines Netz in $z$ ist unnötig. Tabelle~\ref{tab:fem}
fasst die beiden FEM-Modelle zusammen.

\begin{table}[t]
\centering\footnotesize
\caption{Die beiden hochaufgelösten Modelle (gemessen auf einem
Apple-M-Laptop, Python/SciPy).}
\label{tab:fem}
\begin{tabular}{@{}lrr@{}}
\toprule
 & EM-FEM & Wärme-FEM (Scheibe) \\
\midrule
Gleichung & \eqref{eq:em}, komplex & \eqref{eq:heat}, reell, SPD \\
Gebiet & $1\times1\,\mathrm m$ & Scheibe $R{=}80$, $d{=}3\,$mm \\
Knoten / Dreiecke & $8700$ / $17\,028$ & $697$ / $1280$ \\
Lösungszeit & $0{,}40\,$s & $0{,}04\,$s \\
Ergebnis & $q(r,z)$, $P_{\mathrm{Cu}}$, $P_{\mathrm{Fe}}$ & $T(r,z)$ \\
\bottomrule
\end{tabular}
\end{table}

\textbf{Verifikation.} Drei automatische Prüfungen sichern, dass die
Zahlen die Physik und nicht einen Programmierfehler abbilden:
\begin{itemize}[leftmargin=*, nosep]
  \item \emph{Energiebilanz:} erzeugte Wärme $\int q\,dV$ = abgegebene
        Wärme $\oint h(T-T_\infty)\,dA$, Abweichung $0{,}000\,\%$.
  \item \emph{Gebietsgröße:} Dirichlet- gegen Neumann-Rand am
        $1\times1\,\mathrm m$-Gebiet, alle Verluste $<1\,\%$ verschieden
        (größte: $P_{\mathrm{Scheibe}}$, $0{,}77\,\%$).
  \item \emph{Linearität:} volle Neulösung bei doppeltem Strom liefert
        $P(2I)/P(I)=3{,}998$ (Sollwert $4$).
\end{itemize}

\kernidee{Die FEM macht aus einer Feldgleichung ein lineares
Gleichungssystem $\mathbf K\mathbf x=\mathbf f$ mit tausenden
Unbekannten, jede ein Knotenwert. Das ist genau, aber jede Frage („wie
warm bei $6\,$A?“, „wie warm in $30\,$s?“) verlangt eine neue Lösung.}

%%%%%%%%%%%%%%%%%%%%%%%%%%%%%%%%%%%%%%%%%%%%%%%%%%%%%%%%%%%%%%%%%%%%%%%%%%%%%%%%
\section{Warum die FEM allein noch kein Zwilling ist}
\label{sec:why}

Tabelle~\ref{tab:fem} zeigt: In 2D ist die FEM mit $0{,}44\,$s pro
Betriebspunkt schon schnell. Das ehrliche Argument für die Reduktion ist
daher nicht nur Rechenzeit, sondern die Einsatzumgebung eines Zwillings:
\begin{enumerate}[leftmargin=*, nosep]
  \item \textbf{Viele Anfragen statt einer.} Der Zwilling rechnet in
        jedem Zeitschritt, bei jeder Stromänderung, und korrigiert
        $\sigma(T)$ fortlaufend. Mit voller FEM hieße das je
        Schritt eine EM- und eine Wärmelösung.
  \item \textbf{Zielgerät ohne Löser.} Die AR-Ansicht läuft als einzelne
        HTML-Datei im Handy-Browser. Dort gibt es kein SciPy, keinen
        dünnbesetzten LU-Löser und nur wenig Speicher.
  \item \textbf{Kalibrierung.} Aus einer Messkurve lassen sich keine
        $697$ Knotenwerte fitten, wohl aber wenige physikalisch
        interpretierbare Größen wie $\tau$ oder $hA$.
  \item \textbf{Skalierung.} In 3D oder bei gekoppelten Baugruppen wächst
        die FEM auf $10^5$--$10^7$ Unbekannte, dort ist sie für den
        Betrieb schlicht zu langsam~\cite{hartmann2025golden}. Unser 2D-Fall
        ist das überschaubare Lehrbeispiel desselben Prinzips.
\end{enumerate}
Hartmann et al.\ formulieren die Lösung als Trennung in
\emph{Offline}- und \emph{Online}-Phase~\cite{hartmann2018mor}: Offline
wird das große Modell einmal (oder wenige Male) gelöst und daraus ein
kleines Modell extrahiert; online läuft nur noch dieses
(Abb.~\ref{fig:chain}).

%%%%%%%%%%%%%%%%%%%%%%%%%%%%%%%%%%%%%%%%%%%%%%%%%%%%%%%%%%%%%%%%%%%%%%%%%%%%%%%%
\section{Schritt 3: Modellordnungsreduktion}
\label{sec:mor}

\subsection{Die allgemeine Idee}
Das FEM-System \eqref{eq:linsys} hat $N$ Unbekannte. Die meisten davon
bewegen sich aber nicht unabhängig: Heizt man die Scheibe auf, ändert sich
vor allem \emph{wie stark} sie warm ist, kaum \emph{wo}. Projektionsbasierte
MOR~\cite{benner2015survey} nutzt das und schreibt das Feld als
Kombination weniger fester Formen (Basisvektoren) $\mathbf V=[\mathbf
v_1,\dots,\mathbf v_r]$ mit $r\ll N$:
\begin{equation}
  \mathbf T(t)\approx T_\infty+\mathbf V\boldsymbol\beta(t),\qquad
  \underbrace{\mathbf V^{\!\top}\!\mathbf M\mathbf V}_{r\times r}\dot{\boldsymbol\beta}
  +\underbrace{\mathbf V^{\!\top}\!\mathbf K\mathbf V}_{r\times r}\boldsymbol\beta
  =\mathbf V^{\!\top}\tilde{\mathbf f}.
  \label{eq:galerkin}
\end{equation}
Statt $N$ Knotentemperaturen werden nur noch $r$ Amplituden $\beta$
integriert. Die Kunst liegt in der Wahl von $\mathbf V$ und im Nachweis,
dass wenige Formen ausreichen.

\kernidee{Wie ein Foto, das nur heller oder dunkler wird: Statt jedes
Pixel einzeln zu speichern, speichert man das Bild einmal und merkt sich
pro Zeitpunkt nur einen Helligkeitsfaktor.}

Im Projekt kommen drei Reduktionen zum Einsatz, jede durch eine
physikalische Eigenschaft begründet.

\subsection{Reduktion A: Linearität der Elektromagnetik ($I^2$-Skalierung)}
Gl.~\eqref{eq:em} ist bei festem $\omega$ und konstantem $\mu_r$ linear
in $J_s\propto I$, also ist $A_\varphi\propto I$ und nach \eqref{eq:q}
$q\propto I^2$. Die \emph{Form} der Verlustkarte hängt vom Strom gar nicht
ab. Die EM-FEM wird daher einmal beim Referenzstrom $I_{\mathrm{ref}}$
gelöst, danach gilt exakt
\begin{equation}
  q(r,z;I,\bar T)=\hat q(r,z)\Big(\frac{I}{I_{\mathrm{ref}}}\Big)^{\!2}
  \frac{\sigma(\bar T)}{\sigma(T_{\mathrm{ref}})},\quad
  \sigma(T)=\frac{\sigma_0}{1+\alpha(T-T_0)} .
  \label{eq:scale}
\end{equation}
Der Faktor $\sigma(\bar T)$ ($\alpha\approx3{,}9\cdot10^{-3}\,\mathrm K^{-1}$)
berücksichtigt, dass heißes Aluminium schlechter leitet und damit weniger
Wirbelstromverlust erzeugt. Da sich Temperatur über Minuten und das
Feld sofort einstellt, wird $\bar T$ aus dem vorigen Zeitschritt genommen.
In der Sprache der MOR ist das eine parametrische Reduktion mit einem
einzigen, exakt bekannten Basisvektor $\hat q$: $8700$ komplexe
Unbekannte werden zu einer Multiplikation. Gültig ist das, solange das
Eisen nicht sättigt; berechnet wird $B_{\max}=0{,}66\,\mathrm T$
(Effektivwert, Spitze $0{,}94\,\mathrm T$) gegen
$B_{\mathrm{sat}}\approx1{,}5\,\mathrm T$.

\subsection{Reduktion B: Eine Mode für die Scheibe}
Für die Wärme-FEM der Scheibe wählt der Code als einzigen Basisvektor das
stationäre Temperaturfeld beim Referenzstrom,
$\mathbf v_1=\Delta\mathbf T_{\mathrm{ref}}=\mathbf T_{\mathrm{FEM}}-T_\infty$
(ein \emph{Snapshot}). Damit wird \eqref{eq:galerkin} zu einer einzigen
Gleichung erster Ordnung:
\begin{equation}
  \tau\,\dot\beta=\Big(\frac{I}{I_{\mathrm{ref}}}\Big)^{\!2}s(\beta)-\beta,
  \qquad T(r,z,t)=T_\infty+\beta(t)\,\Delta T_{\mathrm{ref}}(r,z),
  \label{eq:rom1}
\end{equation}
mit dem $\sigma(T)$-Faktor $s(\beta)$ aus \eqref{eq:scale} und
$\tau=C/UA$, wobei $C=\rho c_p V=146{,}6\,\mathrm{J/K}$ und
$UA=P_{\mathrm{ref}}/\overline{\Delta T}_{\mathrm{ref}}$.

\textbf{Warum reicht eine Mode?} Das lässt sich prüfen statt behaupten.
Das verallgemeinerte Eigenwertproblem $\mathbf K\boldsymbol\varphi=\lambda
\mathbf M\boldsymbol\varphi$ der Wärme-FEM liefert die Zeitkonstanten
$\tau_i=1/\lambda_i$ aller $697$ Eigenformen. Die ersten lauten
\[
  \tau_1=203{,}9\,\mathrm s,\quad \tau_2=4{,}38\,\mathrm s,\quad
  \tau_3=1{,}33\,\mathrm s,\quad \tau_4=0{,}64\,\mathrm s .
\]
Alle Formen außer der ersten klingen also nach wenigen Sekunden ab, während
die erste die Minuten-Dynamik trägt. Diese \emph{spektrale Lücke}
$\tau_1/\tau_2\approx47$ ist die mathematische Rechtfertigung für $r=1$.
Physikalisch ist es die kleine Biot-Zahl
$\mathrm{Bi}=h\,(d/2)/k\approx25\cdot0{,}0015/237\approx2\cdot10^{-4}$:
Aluminium leitet so gut, dass die Scheibe nahezu gleichmäßig warm wird
($\Delta T_{\mathrm{ref}}$ variiert nur zwischen $42{,}5$ und
$43{,}4\,\mathrm K$).

\begin{figure}[t]
\centering
\includegraphics[width=\columnwidth]{Figures/rom_vs_fem.pdf}
\caption{Aufheizen der Scheibe bei $I_{\mathrm{ref}}$ aus
Umgebungstemperatur. Schwarz: volle instationäre FEM ($697$ Unbekannte,
implizites Euler, $\Delta t=1\,$s). Blau: Ein-Moden-ROM mit
$\tau_1$ aus dem Eigenproblem. Orange: ROM, wie es im Code steht
($\tau=C/UA=244\,$s). Unten: größter Fehler über alle Knoten.}
\label{fig:romfem}
\end{figure}

\textbf{Ergebnis gegen die volle FEM.} Abbildung~\ref{fig:romfem}
vergleicht die reduzierte mit der vollen instationären Lösung. Mit der
Zeitkonstante $\tau_1$ aus dem Eigenproblem ist der größte Fehler über alle
Knoten und Zeiten $0{,}36\,\mathrm K$ ($0{,}8\,\%$ des Endwerts); die
Reduktion $697\to1$ kostet also praktisch keine Genauigkeit. Das im Code
verwendete $\tau=C/UA=244\,\mathrm s$ liegt dagegen $20\,\%$ über $\tau_1$
und erzeugt während des Aufheizens bis zu $3{,}0\,\mathrm K$ ($6{,}9\,\%$)
Abweichung; stationär stimmen beide exakt.

\begin{remark}[Ursache der $20\,\%$]
$UA$ wird gegen $T_\infty=20\,^\circ$C gebildet, die Unterseite der
Scheibe sieht aber die von den Spulen vorgewärmte Luft
($T_{\infty,\mathrm{bot}}=30\,^\circ$C). Das effektive $UA$ ist dadurch
kleiner als die tatsächliche Wärmeabgabe $\mathbf 1^{\!\top}\mathbf
K\mathbf 1=0{,}719\,\mathrm{W/K}$, aus der $\tau_1=C/0{,}719$ folgt. Im
Zwilling wird die Luftvorwärmung allerdings ohnehin als eigener, langsamer
Anteil $\beta_{\mathrm{air}}$ modelliert (Abschnitt~\ref{sec:code}),
während die FEM-Referenz hier eine sofort warme Luft annimmt. Welche
Zeitkonstante am realen Aufbau richtig ist, kann erst eine gemessene
Aufheiz-/Abkühlkurve entscheiden (Abschnitt~\ref{sec:results}).
\end{remark}

\subsection{Reduktion C: Konzentriertes RC-Netz für Spulen und Eisen}
Spulen und Eisen haben keine eigene Wärme-FEM. Sie werden als thermisches
RC-Netz modelliert, wie es für Elektromotoren etabliert
ist~\cite{wallscheid2021thermal}: Jeder Körper ist ein Knoten mit
Wärmekapazität $C_i$, verbunden über Leitwerte $G_{ij}$ und an die Luft
über $hA_i$:
\begin{equation}
  C_i\,\dot T_i=P_i(I)+\sum_j G_{ij}(T_j-T_i)-hA_i\,(T_i-T_{\mathrm{air}}).
  \label{eq:rc}
\end{equation}
Das Netz hat sechs Knoten: je Spule eine IR-sichtbare Oberfläche und ein
träges Wicklungsinneres, den Eisenkern und einen gemeinsamen Luftknoten.
Die Verluste $P_i$ kommen aus der EM-FEM und skalieren nach \eqref{eq:scale};
$C_i$ folgt aus Geometrie und Material. Dieses Modell ist nicht aus der FEM
projiziert, sondern physikalisch aufgestellt (\emph{grey box}). Sein
Vorteil ist Kalibrierbarkeit: Jeder Parameter hat eine Bedeutung und lässt
sich einzeln an Messdaten anpassen (Schritt~4).

\kernidee{Aus $8700+697$ Unbekannten werden acht Zustandsgrößen: zwei
Amplituden für die Scheibe und sechs Knotentemperaturen. Jede Reduktion
ist durch eine prüfbare Eigenschaft begründet: Linearität (A), spektrale
Lücke (B), gut leitende, kompakte Körper (C).}

%%%%%%%%%%%%%%%%%%%%%%%%%%%%%%%%%%%%%%%%%%%%%%%%%%%%%%%%%%%%%%%%%%%%%%%%%%%%%%%%
\section{Schritt 4: Kalibrierung und Betrieb}
\label{sec:calib}

\textbf{Kalibrierung.} Unsicher sind vor allem die Wärmeübergänge
$hA_{\mathrm{innen}}$, $hA_{\mathrm{außen}}$. Sie werden gegen die
IR-Messung im stationären Zustand bei $7{,}8\,\mathrm A$ gefittet
(Innenspule $79\,^\circ$C, Außenspule $74\,^\circ$C, Umgebung $29\,^\circ$C).
\texttt{refit\_hA.py} löst dazu per Nullstellensuche das
\emph{vollständige} nichtlineare Online-Modell bis zum stationären Zustand,
nicht eine von Hand linearisierte Formel. Ergebnis:
$hA_{\mathrm{innen}}=2{,}474\,\mathrm{W/K}$,
$hA_{\mathrm{außen}}=2{,}889\,\mathrm{W/K}$.

\textbf{Betrieb.} Online erhält das Modell jede Sekunde den gemessenen
Effektivstrom (Hall-Sensor ACS712 am Arduino, \texttt{data\_io.py}) und
macht einen Zeitschritt. Ein Schritt ist reine Skalararithmetik und kostet
in Python $\approx3{,}2\,\mu\mathrm s$; dieselbe Rechnung läuft als
JavaScript in der AR-Ansicht auf dem Smartphone. Das gespeicherte
Temperaturfeld $\Delta T_{\mathrm{ref}}(r,z)$ wird mit $\beta(t)$
multipliziert und für die 3D-Darstellung um die Achse rotiert.

%%%%%%%%%%%%%%%%%%%%%%%%%%%%%%%%%%%%%%%%%%%%%%%%%%%%%%%%%%%%%%%%%%%%%%%%%%%%%%%%
\section{Code-Dokumentation}
\label{sec:code}

Dieser Abschnitt beschreibt, was der Code tut, sodass er ohne Lektüre des
Quelltexts nachvollziehbar ist. Tabelle~\ref{tab:modules} ordnet jede Datei
einem Kasten aus Abb.~\ref{fig:chain} zu.

\begin{table*}[t]
\centering\footnotesize
\caption{Module, ihre Aufgabe, Ein-/Ausgaben und Methode. Reihenfolge = Datenfluss.}
\label{tab:modules}
\begin{tabularx}{\textwidth}{@{}l c X X X@{}}
\toprule
Datei & Schritt & Aufgabe & Eingang $\to$ Ausgang & Methode / Bibliothek \\
\midrule
\texttt{params.yaml} & 1 & Einzige Quelle aller Zahlen (Geometrie, Material, Randbedingungen, Netz, Kalibrierwerte) & -- & YAML; keine Konstante im Code \\
\texttt{config.py} & 1 & Lädt und normiert Parameter (mm$\to$m, $I_{\mathrm{peak}}=\sqrt2I_{\mathrm{rms}}$) & YAML $\to$ \texttt{Config}-Objekt & PyYAML \\
\texttt{em\_solver.py} & 2 & EM-Phasor-FEM Gl.~\eqref{eq:em}; Verluste, Hubkraft, Sättigungs- und Gebietsprüfung & Geometrie, $I$ $\to$ $A_\varphi$, $q(r,z)$, $P_{\mathrm{Scheibe}}, P_{\mathrm{Fe}}, P_{\mathrm{Cu}}$, $F_z$ & P1-Galerkin, komplex, \texttt{spsolve} \\
\texttt{thermal\_solver.py} & 2 & Stationäre Wärme-FEM Gl.~\eqref{eq:weak} der Scheibe; Energiebilanz & $q(r,z)$ $\to$ $T(r,z)$ & P1-Galerkin, \texttt{spsolve}, Interpolation EM$\to$Wärmenetz \\
\texttt{rom.py} & 3 & Baut Reduktion A+B: speichert $\Delta T_{\mathrm{ref}}$, berechnet $\tau$, $UA$ & FEM-Ergebnis $\to$ ROM-Koeffizienten & Snapshot, Energiebilanz \\
\texttt{twin\_model.py} & 3 & Brücke: berechnet alle Koeffizienten (inkl. RC-Netz) je Scheibenradius & Config, EM, ROM $\to$ Koeffizientensätze & Cache je Scheibe \\
\texttt{twin\_core.py} & 3 & \textbf{Der Online-Zwilling.} \texttt{TwinState.step(I,\,dt)} integriert Gl.~\eqref{eq:rom1} + \eqref{eq:rc} + Levitation & $I(t)$ $\to$ $\beta$, $T_i$, $z_{\mathrm{gap}}$ & expl. Euler mit Teilschritten $\Delta t\le0{,}05\tau$; nur NumPy + stdlib \\
\texttt{refit\_hA.py} & 4 & Kalibriert $hA_{\mathrm{innen/außen}}$ durch das echte Online-Modell & IR-Messwerte $\to$ \texttt{params.yaml} & Nullstellensuche \\
\texttt{data\_io.py} + Arduino & 4 & Liest Messstrom (1\,Hz, CSV über USB), treibt \texttt{TwinState} & Seriell $\to$ $I_{\mathrm{rms}}(t)$ & RMS über 5 Netzperioden \\
\texttt{build\_twin\_html\_fem.py} & 4 & „Backt“ Koeffizienten + Online-Modell als JavaScript in eine HTML-Datei (3D/AR) & Koeffizienten $\to$ \texttt{digital\_twin\_fem.html} & prozedurale Geometrie, three.js \\
\texttt{xval\_twin.py} & -- & Prüft, dass Python- und JS-Zwilling identisch rechnen & beide $\to$ PASS/FAIL & Playwright, Toleranz $10^{-9}$ \\
\bottomrule
\end{tabularx}
\end{table*}

\textbf{Ein Durchlauf.} Ruft man \texttt{python build\_twin\_html\_fem.py}
auf, geschieht Folgendes: (1) \texttt{config.py} liest
\texttt{params.yaml}. (2) \texttt{em\_solver.py} baut das graduierte
Dreiecksnetz des $1\times1\,$m-Gebiets, weist jedem Element $\nu$,
$\sigma$ und $J_s$ zu, assembliert und löst das komplexe System und
integriert daraus die Verluste je Körper. (3) \texttt{thermal\_solver.py}
interpoliert $q$ auf das Scheibennetz, löst die Wärmeleitung und prüft die
Energiebilanz. (4) \texttt{rom.py} speichert das Feld als Mode und
berechnet $\tau$. Zusätzlich wird aus einem zweiten Lauf ohne
Spulenluftvorwärmung der Anteil $f_{\mathrm{eddy}}$ bestimmt, den der
Wirbelstrom am Scheibenfeld hat ($\approx0{,}84$); der Rest
$f_{\mathrm{air}}\approx0{,}16$ folgt der Luft um die Spulen. Online ist
$\beta=f_{\mathrm{eddy}}\beta_{\mathrm{eddy}}+f_{\mathrm{air}}\beta_{\mathrm{air}}$.
(5) Die Koeffizienten des RC-Netzes werden aus Geometrie und
Kalibrierwerten berechnet. (6) Alles wird zusammen mit dem Online-Modell
als JavaScript in eine einzige HTML-Datei geschrieben, die ohne eigenen
Server im Browser läuft (nur three.js wird vom CDN geladen).

\textbf{Zwei Implementierungen, ein Modell.} Das Online-Modell existiert
zweimal: als \texttt{twin\_core.py} (Python-Referenz, genutzt von allen
Desktop-Ansichten und dem Sensorpfad) und als eingebettetes JavaScript.
\texttt{xval\_twin.py} startet die HTML-Datei in einem Browser und vergleicht
beide über sieben Stromverläufe (absolute Toleranz $10^{-9}$). Eine zweite
Prüfung stellt fest, ob die eingebackenen Koeffizienten noch zu
\texttt{params.yaml} passen. Beide müssen nach jeder Änderung bestehen.

\textbf{Levitation.} Nebenbei berechnet der Zwilling die Schwebehöhe: Die
EM-FEM liefert die Hubkraft $F_z(z)$, deren Gleichgewicht mit der
Gewichtskraft $z_{\mathrm{eq}}$ ergibt; die Bewegung dorthin wird als
gedämpfter Feder-Masse-Schwinger in geschlossener Form berechnet.

%%%%%%%%%%%%%%%%%%%%%%%%%%%%%%%%%%%%%%%%%%%%%%%%%%%%%%%%%%%%%%%%%%%%%%%%%%%%%%%%
\section{Endergebnisse}
\label{sec:results}

\begin{table}[t]
\centering\footnotesize
\caption{Verlustleistung aus der EM-FEM ($I=5\,\mathrm A$, $T_\infty=20\,^\circ$C).}
\label{tab:losses}
\begin{tabular}{@{}lr@{}}
\toprule
Quelle & $P$ [W] \\
\midrule
Wirbelstrom Scheibe (Al) & 25{,}81 \\
Wirbelstrom Eisenkern + -ring & 5{,}86 \\
Ohmsch Spulen (innen 52{,}32 + außen 54{,}12) & 106{,}44 \\
\midrule
Gesamt & 138{,}11 \\
\bottomrule
\end{tabular}
\end{table}

\begin{table}[t]
\centering\footnotesize
\caption{Aufwand: volles Modell gegen reduziertes Modell.}
\label{tab:speed}
\begin{tabular}{@{}lrr@{}}
\toprule
 & FEM (offline) & ROM (online) \\
\midrule
Unbekannte & $8700$ (kompl.) + $697$ & $8$ Zustände \\
Aufwand & $0{,}44\,$s / Betriebspunkt & $3{,}2\,\mu$s / Zeitschritt \\
Benötigt & SciPy, dünnbes. LU & Arithmetik \\
Läuft auf & Laptop & Laptop, Browser, Handy \\
Fehler Scheibe & Referenz & $\le0{,}36\,$K (1 Mode, $\tau_1$) \\
\bottomrule
\end{tabular}
\end{table}

\textbf{Verluste.} Tabelle~\ref{tab:losses}: Die Spulen erzeugen $77\,\%$
der Wärme, die Scheibe $19\,\%$. Das bestätigt die IR-Beobachtung, dass
die Innenspule das heißeste Bauteil ist.

\textbf{Reduktion.} Tabelle~\ref{tab:speed}: Das Online-Modell ist um
rund fünf Größenordnungen schneller als eine volle Neulösung und kommt ohne
Löserbibliothek aus. Die Eigenwertanalyse belegt, dass die Reduktion der
Scheibe auf eine Mode praktisch verlustfrei ist (Abb.~\ref{fig:romfem}).

\textbf{Stationäre Validierung.} Am Kalibrierpunkt ($7{,}8\,\mathrm A$)
trifft der Zwilling die gemessenen Spulentemperaturen exakt
($79{,}00/74{,}00\,^\circ$C). Das ist ein Kalibrier-, kein
unabhängiger Validierungsnachweis.

\textbf{Transiente: nicht validiert.} Gegen die einzige vorhandene
zeitaufgelöste Messung (Stromrampe, fünf IR-Zeitpunkte der Außenspule)
heizt das Modell systematisch zu langsam auf (RMS-Fehler
$11{,}5\,\mathrm K$). Kapazitätsaufteilung und innere Leitwerte des RC-Netzes
sind aus reinen Aufheizdaten nicht eindeutig bestimmbar; dafür fehlt eine
gemessene Abkühlkurve.

\textbf{Schwebehöhe: offen.} Berechnet $z_{\mathrm{eq}}\approx11{,}7\,$mm
(Unterkante), beobachtet ist ein sichtbarer Spalt von $7$--$8\,$mm.
Sättigung ist als Ursache ausgeschlossen ($<0{,}1\,\%$ Kraftänderung);
wahrscheinlich ist der nicht gemessene Wert $\mu_r=1000$ des Eisens.
$\mu_r$ wurde bewusst nicht auf die Beobachtung zurückgefittet.

\begin{remark}[Stromkonvention]
$5\,\mathrm A$ ist ein gemessener Effektivwert. Die Verlustkette nutzt ihn
bewusst als Amplitude; der konstante Faktor $2$ in der Leistung wird beim
Fit von $hA$ absorbiert, die Temperaturen bleiben korrekt. Die Kraftkette
rechnet mit der echten Amplitude $\sqrt2\cdot5=7{,}07\,\mathrm A$, weil
sich eine Kraft nicht wegkalibrieren lässt.
\end{remark}

%%%%%%%%%%%%%%%%%%%%%%%%%%%%%%%%%%%%%%%%%%%%%%%%%%%%%%%%%%%%%%%%%%%%%%%%%%%%%%%%
\section{Lessons Learned}
\label{sec:lessons}

\textbf{Was gut lief.}
\begin{itemize}[leftmargin=*, nosep]
  \item \emph{Physik zuerst reduzieren.} Achsensymmetrie (3D$\to$2D),
        Phasor ($50\,$Hz nicht integrieren) und Linearität ($I^2$)
        haben den größten Teil des Aufwands eingespart, bevor überhaupt
        MOR-Technik nötig war.
  \item \emph{Harte, automatische Prüfungen.} Energiebilanz,
        Gebietsgröße, $I^2$-Test und der Python/JS-Abgleich laufen nach
        jeder Änderung und haben mehrere Fehler gefunden, die sonst
        unbemerkt geblieben wären.
  \item \emph{Eine Datei für alle Zahlen.} \texttt{params.yaml} als
        einzige Quelle machte die Nachvermessung der Geometrie zu einer
        reinen Dateiänderung.
  \item \emph{Interpretierbares ROM.} Jeder Online-Parameter
        ($\tau$, $hA$, $C$) ist eine physikalische Größe und damit an
        Messungen nachstellbar.
\end{itemize}

\textbf{Was schlecht lief.}
\begin{itemize}[leftmargin=*, nosep]
  \item \emph{Zu wenig Messdaten, zu spät.} Ohne Abkühlkurve bleibt die
        Transiente unvalidiert; die Sensorhardware hätte früher
        geplant werden müssen.
  \item \emph{Effektiv- und Spitzenwert verwechselt.} Der Fehler
        $I_{\mathrm{rms}}$ vs.\ $\hat I$ wurde zweimal an übersehenen
        Aufrufstellen ausgeliefert. Lehre: physikalische Konventionen im
        Code benennen (\texttt{I\_peak}), nicht im Kopf behalten.
  \item \emph{Kalibrierung außerhalb des Modells.} Eine von Hand
        linearisierte Fit-Formel passte nicht zum nichtlinearen
        Online-Modell (Spulen $7\,$K zu kalt). Lehre: immer durch das
        Modell fitten, das tatsächlich läuft.
  \item \emph{ROM nicht direkt gegen sein FEM geprüft.} Die Abweichung
        $\tau=244\,$s vs.\ $\tau_1=204\,$s fiel erst beim Schreiben
        dieses Berichts durch die Eigenwertanalyse auf. Ein
        ROM-gegen-FEM-Test gehört von Anfang an in die Prüfkette.
  \item \emph{Geometrie spät vermessen.} Bis Juli wurde mit geschätzten
        Radien gerechnet; nach der Nachmessung stiegen die Scheibenverluste
        von $3{,}0$ auf $25{,}8\,$W.
\end{itemize}

%%%%%%%%%%%%%%%%%%%%%%%%%%%%%%%%%%%%%%%%%%%%%%%%%%%%%%%%%%%%%%%%%%%%%%%%%%%%%%%%
\section{Zusammenfassung und Ausblick}
\label{sec:summary}

Der thermische Zwilling folgt der Offline/Online-Trennung der
Modellordnungsreduktion~\cite{hartmann2018mor}: Zwei FEM-Modelle werden
einmal gelöst, drei physikalisch begründete Reduktionen machen daraus ein
Modell mit acht Zuständen, das im Mikrosekundentakt auf einem Smartphone
läuft und aus dem gemessenen Strom die Temperatur vorhersagt. Die
Eigenwertanalyse zeigt, dass diese Reduktion für die Scheibe nahezu
verlustfrei ist, und hat zugleich eine $20\,\%$-Abweichung in der
verwendeten Zeitkonstante aufgedeckt. Nächste Schritte: (1) eine gemessene
Aufheiz- und Abkühlkurve, um $\tau$ und das RC-Netz zu identifizieren;
(2) $\tau$ der Scheibe aus dem Eigenproblem statt aus $C/UA$ ableiten,
sobald die Messung zeigt, welche Definition den Aufbau besser trifft;
(3) eine Messung von $\mu_r$ zur Klärung der Schwebehöhe.

\bibliographystyle{IEEEtran}
\bibliography{dt4tm_references}

\end{document}
